
%%FILE blueprint/src/content.tex
\chapter{Introduction}
The goal of this project is to formalize Schlessinger's criterion in Lean following
the original article \cite{Sch}. Some parts are also based on the proof given in
\cite[\href{https://stacks.math.columbia.edu/tag/06G7}{Tag 06G7}]{stacks-project}.


\chapter{Definitions}
\label{sec:defs}

\section{Base category}

Let $\Lambda$ be a complete Noetherian (commutative) local ring and $k$ its
residue field.

\begin{definition}
  \label{def:ArtinResAlg}
  \lean{Deformation.ArtinResAlg}
  \leanok
  We define $\CL$ to be the category of Artinian local $\Lambda$-algebras
  having residue field $k$, where morphisms are local algebra homomorphisms.
\end{definition}

The condition of having the same residue field for a $\Lambda$-algebra $A$ can
be stated
in a more general context as follows: the algebra map $\Lambda \to \rf_A$
is surjective (where $\rf_A$ is the residue field of $A$).
As such the conditions on $\Lambda$ are omitted in Lean.

Furthermore, we do not need to assume the morphisms to be local, as we have the
following result:

\begin{lemma}
  \label{lem:islocalhom_of_artinian}
  \lean{Deformation.instIsLocalHomAlgHomOfIsArtinianRing_schlessinger}
  \leanok
  Let $A$ and $B$ be commutative local rings, and $f : A \to B$ a ring morphism.
  If $A$ is Artinian, then $f$ is local.
\end{lemma}

\begin{proof}
  \leanok
  In an Artinian local ring, the maximal ideal is the set of
  nilpotent elements. In particular if $x \in \m_A$, then $f (x)$
  is not a unit.
\end{proof}

The Artinian condition on elements of $\CL$ could potentially be
interpreted in two ways: either as Artinian rings or as Artinian
$\Lambda$-modules. In this context, both are equivalent.
\begin{lemma}
  \label{lem:lengtheq}
  \lean{Deformation.length_eq}
  \leanok
  Let $A$ be a local $\Lambda$-algebra with same residue field as $\Lambda$.
  Then $\length_\Lambda (A) = \length_A (A)$.
\end{lemma}

\begin{proof}
  \leanok
  We prove a more general result: if $M$ is an $A$-module,
  then $\length_\Lambda (M) = \length_A (M)$.
  This result follows easily by induction on $\length_\Lambda (M)$ from the
  case $\length_\Lambda (M) = 1$, in which case we have $M \cong k$.

  \cite[\href{https://stacks.math.columbia.edu/tag/06GG}{Tag 06GG}]{stacks-project}
\end{proof}

\begin{lemma}
  \label{lem:has_pullbacks}
  \lean{instHasPullbacksArtinResAlg}
  \leanok
  \uses{def:ArtinResAlg}
  The category $\CL$ has pullbacks.
\end{lemma}

\begin{proof}
  \leanok
  \uses{lem:lengtheq}
  Let $f : X \to Z$ and $g : Y \to Z$ be morphisms in $\CL$.
  Consider the set $P = \{(x, y) \in X \times Y ~\vert~ f (x) = g (y)\}$.
  Since $f$ and $g$ are local, the subset $\m_X \times \m_Y$ is an ideal
  (which is maximal).
  We have $\rf_\Lambda \to \rf_P \to \rf_X$, so $\rf_P \cong k$.
  Furthermore, since $X$ and $Y$ are Artinian $\Lambda$-modules
  by \cref{lem:lengtheq}, $P$
  is an Artinian $\Lambda$-module as well, so it is an Artinian ring.
\end{proof}


We also note that $k$ is a terminal object in $\CL$, so the category also
admits products (which are the same as pullbacks over the terminal object).

\section{Small Extensions}

Let $p : B \to A$ be a surjective morphism in $\CL$.

\begin{definition}
  \label{def:IsSmallExtension}
  \lean{IsSmallExtension}
  \leanok
  \uses{def:ArtinResAlg}
  $p$ is a \emph{small extension} if $\ker p$ is a non-zero principal
  ideal such that $\m_B \cdot \ker p = 0$.
\end{definition}

\begin{lemma}
  \label{lem:smallext_iff}
  \lean{IsSmallExtension.isSmallExtension_iff, IsSmallExtension.isSmallExtension_iff_length_ker}
  \leanok
  \uses{def:IsSmallExtension}
  The following propositions are equivalent:
  \begin{enumerate}
    \item
    $p$ is a small extension
    \item
    $\ker p$ is a minimal non-zero ideal
    \item
    $\length_\Lambda (B) = \length_\Lambda (A) + 1$
  \end{enumerate}
\end{lemma}

\begin{proof}
  \leanok
  Since $p$ is surjective, we have
  $\length_\Lambda (B) = \length_\Lambda (A) + \length_\Lambda (\ker p)$,
  and we also have that $\length_B (\ker p) = 1$ if and only if
  $\ker p$ is a minimal non-zero ideal,
  so using \cref{lem:lengtheq} we get $(2) \iff (3)$.
  We now show $(1) \iff (2)$.

  Assume that $p$ is a small extension, write $\ker p = (t)$,
  let $I < \ker p$ be an ideal and $x \in I$.
  There exists $y \in B$ such that $x = ty$. Since $I \ne (t)$,
  $y$ is not a unit, so we have $y \in \m_B$, and $x = 0$
  by definition of a small extension.
  So $I=0$ and $\ker p$ is minimal.

  Assume that $\ker p$ is a minimal non-zero ideal. Let
  $t \in \ker p \setminus \{0\}$. By minimality, we have $\ker p = (t)$.
  Furthermore, since  $\m_B$ is nilpotent and $\ker p \neq 0$ we have
  $\m_B \cdot \ker p \neq \ker p$ so $\m_B \cdot \ker p = 0$ by minimality.
  So $p$ is a small extension.
\end{proof}

\begin{lemma}
  \label{lem:smallext_hom_induction}
  \lean{IsSmallExtension.hom_induction}
  \leanok
  \uses{def:IsSmallExtension}
  If a property $C$ on morphisms of $\CL$ is satisfied by isomorphisms
  and stable by composition on the left by a small extension, then
  it holds for any surjection.
\end{lemma}

\begin{proof}
  \leanok
  \uses{lem:smallext_iff}
  By induction on $\length_\Lambda (B) \in \N$.
  If $\ker p = 0$, then $p$ is an isomorphism.
  Otherwise, there exists a minimal non-zero ideal $I$ of $B$ because $B$ is
  Artinian. By minimality, since $\ker p \neq 0$,
  we have $I \le \ker p$, so $p$ factorizes through
  the small extension $B \to B/I$.
  We conclude by applying the induction
  hypothesis to the map $B/I \to A$.
\end{proof}

\begin{lemma}
  \label{lem:smallext_obj_induction}
  \lean{IsSmallExtension.obj_induction}
  \leanok
  \uses{def:IsSmallExtension}
  If a property $C$ on objects of $\CL$ is satisfied by $k$,
  stable by isomorphisms
  and stable by small extensions, then
  it holds for all objects of $\CL$.
\end{lemma}

\begin{proof}
  \leanok
  One way is to apply the previous lemma to the map $p : B \to k$
  (i.e. with the property on morphisms $q : X \to Y$ given by if $Y = k$ then $C (X)$).
  In the code the proof is essentially copied and adapted from the previous lemma.
\end{proof}

\section{Completed base category}

For pro-representability, we consider a bigger category with objects
that are projective limits of objects in the base category.
However, in this context the term pro-representability is used in a
more restrictive sense than usual, so that the "completed" category
is not the pro-category as one might expect.

\begin{definition}
  \label{def:complete}
  \lean{Deformation.IsComplete}
  \leanok
  A local ring with maximal ideal $\m$
  is \emph{complete} if it is $\m$-adically complete.
\end{definition}

\begin{definition}
  \label{def:ProArtinResAlg}
  \lean{Deformation.ProArtinResAlg}
  \leanok
  \uses{def:complete}
  We define $\Chat$ to be the category of complete Noetherian local
  $\Lambda$-algebras having residue field $k$,
  where morphisms are local algebra homomorphisms.
\end{definition}

Note that a local ring is endowed with a natural topological ring structure,
namely the $\m$-adic topology, where $\m$ is the maximal ideal.
It is defined by setting the powers of $\m$
as a neighborhood base of $0$.

\begin{lemma}
  \label{lem:fullsubcategory}
  \lean{Deformation.inclusionFunctor}
  \leanok
  \uses{def:ProArtinResAlg, def:ArtinResAlg}
  The category $\Chat$ is a full subcategory of $\CL$.
\end{lemma}

\begin{proof}
  \leanok
  An Artinian ring is Noetherian, and its $\m$-adic topology is discrete,
  so it is trivially complete.
\end{proof}

\begin{lemma}
  \label{lem:isArtinian_of_noetherian_nilpotent}
  \lean{Deformation.isArtinian_of_noetherian_nilpotent}
  \leanok
  A local Noetherian ring whose maximal ideal is nilpotent is Artinian.
\end{lemma}

\begin{proof}
  \leanok
  A Noetherian ring is Artinian if and only if it has Krull dimension
  0, so it suffices to show that any prime ideal is maximal. Let $I$
  be a prime ideal. Since $\m^n = 0 \le I$ we have $\m \le I$, so
  $I = \m$ is maximal.
\end{proof}

\begin{lemma}
  \label{lem:quotinC}
  \lean{quot_C, instQuotInCHPowIdealMaximalIdealHAddOfNat}
  \leanok
  \uses{def:ProArtinResAlg, def:ArtinResAlg}
  If $A \in \Chat$, then for any $n \in \N^*$ we have
  $A / \m^n \in \CL$ where $\m$ is the maximal ideal of $A$.
  In fact, we have $A / J \in \CL$ for any ideal $J$ such that
  $\m^n \le J < A$.
\end{lemma}

\begin{proof}
  \leanok
  \uses{lem:isArtinian_of_noetherian_nilpotent}
  $A / \m^n$ is a local Noetherian ring, with nilpotent maximal ideal
  $\m / \m^n$ and residue field $k$, so \cref{lem:isArtinian_of_noetherian_nilpotent}
  applies.
\end{proof}

\begin{remark}
  The fact that $A \in \Chat$ is complete (and Noetherian) means that we have
  $A = \varprojlim A / \m^n$. We don't really use this result here but
  it explains some of the underlying ideas.
\end{remark}

\begin{lemma}
  \label{lem:factorizes}
  \lean{Deformation.maxideal_pow_le_ker_of_artinian}
  \leanok
  Let $R \in \Chat$, $A \in \CL$ and $u : R \to A$. There exists
  $n \in \N^*$ such that $\m_R^n \le \ker u$, so $u$ factorizes
  through $u_n : R/\m_R^n \to A$.
\end{lemma}

\begin{proof}
  \leanok
  There exists $n$ such that $\m_A^n = 0$,
  so $\m_R^n \le u^{-1} (\m_A ^n) = \ker u$.
\end{proof}

\section{Trivial Square-Zero Extensions}
\label{subsubsec:tsze}

\begin{definition}
  \label{def:tsze_C}
  \lean{Deformation.tsze_C}
  \leanok
  \uses{def:ArtinResAlg}
  For $V$ a $k$-vector space, we write $k[V]$ for the
  \emph{trivial square-zero extension}
  of $V$ over $k$, that is the module $k \oplus V$ with
  multiplication defined so that $V ^ 2 = 0$.
\end{definition}

\begin{definition}
  \label{def:dualnumber_C}
  \leanok
  \lean{dualNumber_C}
  \uses{def:tsze_C}
  We write $k[\varepsilon]$ for \emph{dual numbers}, the special case
  where $V = k$.
\end{definition}

\begin{definition}
  \label{def:tsze_functor}
  \leanok
  \lean{Deformation.tsze_functor}
  \uses{def:tsze_C}
  We can a functor $\Mod_k^{fg} \to \CL$
  that commutes with finite products (where $\Mod_k^{fg}$ is the category
  of finite dimensional $k$-vector spaces).
\end{definition}

\begin{lemma}
  \label{lem:vstsze}
  \lean{Deformation.instModuleResidueFieldObjArtinResAlgTsze_C}
  \leanok
  \uses{def:tsze_functor}
  If $F$ is a functor $\CL \to Sets$ such that
  \[ F(k[V] \times k[W]) ~ \tilde{\to} ~ F(k[V]) \times F(k[W]) \]
  for all finite dimensional $k$-vector spaces $V$ and $W$, then
  $F(k[V])$ has a canonical $k$-vector space structure.
\end{lemma}

\begin{proof}
  \leanok
  \proves{lem:vstsze}
  \uses{lem:vs}
  This is a special case of the following lemma
  \cite[\href{https://stacks.math.columbia.edu/tag/06I6}{Tag 06I6}]{stacks-project}:
\end{proof}

\begin{lemma}
  \label{lem:vs}
  \lean{vs.instModuleObjFGModuleCat_schlessinger}
  \leanok
  Let $R$ be a (commutative) ring, and $L : \Mod_R^{fg} \to Sets$ be a functor
  that preserves finite products.
  Then for any $M \in \Mod_R^{fg}$, $L(M)$ has a canonical $R$-module
  structure.
\end{lemma}

\begin{proof}
  \leanok
  Let $M \in \Mod_R^{fg}$. Write $s : M \times M \to M$ for the addition
  map, $\lambda_r : M \to M$ for the scalar multiplication by $r \in R$,
  and $z : 0 \to M$.
  The module structure on $L(M)$ is defined using the images of these
  morphisms by $L$, together with the identification
  $L(M\times M) \cong L(M) \times L(M)$, and the fact that $L(0)$ is terminal
  in $Sets$ so it is a singleton.
  Checking that this defines a module structure is rather straightforward,
  by writing suitable commutative diagrams in $\Mod_R^{fg}$ and
  pushing them with $L$.
\end{proof}

\begin{lemma}
  \label{lem:vsnatlin}
  \lean{vs.linmap_of_nattrans}
  \leanok
  \uses{lem:vs}
  Let $R$ be a (commutative) ring, $L$ and $L' : \Mod_R^{fg} \to Sets$ be two
  functors that preserve finite products, and $\eta : L \to L'$ a natural
  transformation.
  Then for any $M \in \Mod_R^{fg}$, $\eta_M$ is $R$-linear.
\end{lemma}

\begin{proof}
  \leanok
  This is a very simple check.
\end{proof}

\begin{lemma}
  \label{lem:vsnatiso}
  \lean{vs.isIso_natTrans_of_isIso_dim1}
  \leanok
  Let $R$ be a (commutative) ring, $L$ and $L' : \Mod_R^{fg} \to Sets$ be two
  functors that preserve finite products, and $\eta : L \to L'$ a natural
  transformation.
  Then $\eta$ is an isomorphism if and only if
  $\eta_R$ is an isomorphism.
\end{lemma}

\begin{proof}
  \uses{lem:vsnatlin}
  By induction on the dimension.
\end{proof}

\section{Artin Functors}

\begin{definition}
  \label{def:hR}
  \lean{hR}
  \leanok
  \uses{def:ProArtinResAlg, def:ArtinResAlg}
  For $R \in \Chat$ and $A \in \CL$, we define a functor $\CL \to Sets$ by
  $h_R (A) = \Hom_{\Chat}(R, A)$.
\end{definition}

In what follows we will only consider functors $F$ from $\CL$ to $Sets$
that preserve terminal objects,
or equivalently such that $F(k)$ is a singleton.

\begin{definition}
  \label{def:completion}
  \lean{CategoryTheory.Functor.completion}
  \leanok
  \uses{def:ProArtinResAlg, def:ArtinResAlg}
  We can extend such a functor to $\Chat$ by setting for $R \in \Chat$,
  $\hat{F}(R) = \varprojlim F(R / \m^n)$. A map $u : R \to S$ induces maps
  $u_n : R / \m_R^n \to S / \m_S^n$ for all $n \in \N$, so we can also define
  $\hat{F}(u)$ by using the universal property.
  This defines a functor $\hat{F} : \Chat \to Sets$.
\end{definition}

\begin{lemma}
  \lean{natiso_completion}
  \leanok
  \label{lem:isocompl}
  \uses{def:completion}
  There is a natural isomorphism $\hat{F}\restr{\CL} \cong F$.
\end{lemma}

\begin{proof}
  \leanok
  If $A$ is in $\CL$, there exists $N \in \N^*$ such that $\m_A^N=0$, so we
  have $F(A/\m_A^n) \cong F(A)$ for all $n \ge N$.
\end{proof}

\begin{definition}
  \label{def:procouple}
  \lean{Procouple}
  \leanok
  \uses{def:completion}
  A \emph{pro-couple} is a pair $(R, \xi)$ where $R \in \Chat$ and $\xi \in \hat{F} (R)$.
\end{definition}

\begin{definition}
  \lean{nattrans_of}
  \leanok
  \label{def:nattrans_of}
  \uses{def:hR, def:procouple}
  A pro-couple $(R, \xi)$ defines a natural transformation
  $\nattrans(\xi) : h_R \to F$ given by
  $u \mapsto \hat{F} (u) (\xi)$ up to the isomorphism given in
  lemma~\ref{lem:isocompl}.
\end{definition}

\begin{lemma}
  \label{lem:nattrans_of_comp_artinhom}
  \lean{nattrans_of_comp_artinhom}
  \leanok
  \uses{def:nattrans_of}
  Let $(R, \xi)$ be a pro-couple, $A, B \in \CL$, $u : R \to A$ and $v : A \to B$.
  Then $\nattrans(\xi) (v \circ u) = F(v)(\nattrans(\xi) (u))$.
\end{lemma}

\begin{proof}
  \leanok
  This is essentially the naturality of $\hat{F}\restr{\CL} \cong F$.
\end{proof}

\begin{lemma}
  \label{lem:completion_map_quot_mk_eq_pi}
  \lean{completion_map_quot_mk_eq_pi}
  \leanok
  \uses{def:nattrans_of}
  Let $(R, \xi)$ be a pro-couple, with $\xi = (\xi_q)_{q\in\N}$, $n\in\N^*$,
  and $p:R\to R/\m_R^n$ the quotient map. Then $\nattrans(\xi) (p) = \xi_n$.
\end{lemma}

\begin{proof}
  \leanok
  In fact, we have that $\hat{F}(p)$ is equal to the projection
  $\hat{F}(R) \to F(R / \m_R^n)$ composed with the natural isomorphism
  $F \cong \hat{F}\restr{\CL}$.
  This is a simple check. %RK: using the universal property of the limit.
\end{proof}

In particular, since any map $R \to A$ with $A \in \CL$ factorizes through
some $u_n : R / \m^n \to A$ (\cref{lem:factorizes}),
we have $\nattrans(\xi) (u) = F(u_n) (\xi_n)$.

\begin{definition}
  \label{def:Prorepresents}
  \lean{Prorepresents}
  \leanok
  \uses{def:nattrans_of, def:hR}
  A pro-couple $(R, \xi)$ \emph{pro-represents} $F$
  if the natural transformation $\nattrans(\xi) : h_R \to F$ it
  induces is an isomorphism.
  $F$ is \emph{pro-representable} if there exists a pro-couple which
  pro-represents $F$.
\end{definition}

\begin{definition}
  \label{def:IsSmooth}
  \lean{IsSmooth}
  \leanok
  \uses{def:ArtinResAlg}
  Let $F, G : \CL \to Sets$ be two functors that preserve terminal objects.
  A natural transformation $\eta : F \to G$ is \emph{smooth}
  if for any surjection
  $p : B \to A$ in $\CL$,
  the morphism
  \[ F(B) \to F(A) \times_{G(A)} G(B)\]
  is surjective.
\end{definition}

\begin{lemma}
  \label{lem:smoothsmallext}
  \lean{isSmooth_iff}
  \leanok
  \uses{def:IsSmooth, def:IsSmallExtension}
  To check that a natural transformation is smooth, it is enough
  to check the condition on morphisms $p$ that are small extensions.
\end{lemma}

\begin{proof}
  \leanok
  \uses{lem:smallext_hom_induction}
  Apply \cref{lem:smallext_hom_induction}.
  Check that the property
  is stable by composition (note that the pullback is in $Sets$).
\end{proof}

\begin{lemma}
  \label{lem:surj_of_smooth}
  \lean{surj_of_smooth}
  \leanok
  \uses{def:IsSmooth}
  Let $F, G : \CL \to Sets$ be two functors that preserve terminal objects.
  If $\eta : F \to G$ is smooth,
  then $\eta_A$ is surjective for any $A \in \CL$.
\end{lemma}

\begin{proof}
  \leanok
  Apply the condition to $p : A \to k$ (using that $F(k)$ and $G(k)$
  are singletons).
\end{proof}

\begin{definition}
  \label{def:tangentSpace}
  \lean{tangentSpace}
  \leanok
  \uses{def:dualnumber_C}
  The \emph{tangent space} of a functor $F$ is $t_F = F(k[\varepsilon])$.
  In the case $F = h_R$, we simply write $t_R$.
\end{definition}

\begin{definition}
  \label{def:IsHull}
  \lean{IsHull}
  \leanok
  \uses{def:nattrans_of, def:tangentSpace, def:IsSmooth}
  A pro-couple $(R, \xi)$ is a \emph{hull} of $F$ if the natural
  transformation
  $\nattrans(\xi) : h_R \to F$ it induces
  is smooth and $\nattrans(\xi)_{k[\varepsilon]} : t_R \to t_F$
  is an isomorphism.
\end{definition}

Note that if $\eta : F \to G$ is an isomorphism on $k[\varepsilon]$,
then it is an isomorphism on $k[V]$ for any finite dimensional $k$-vector space
(lemma~\ref{lem:vsnatiso}).

\section{Schlessinger's conditions}

Let $F$ be a functor from $\CL$ to $Sets$ such that $F(k)$ is a singleton.
For $f : X \to Z$ and $G : Y \to Z$ morphisms in
 $\CL$, we consider the comparison morphism:
\begin{equation}
  \label{eq:pfg}
  p_{f, g} : F(X \times_Z Y) \to F(X) \times_{F(Z)} F(Y)
\end{equation}

\begin{definition}
  \label{def:H1}
  \lean{H1}
  \leanok
  \uses{def:IsSmallExtension, lem:has_pullbacks}
  $(H_1)$ : if $g$ is a small extension then $p_{f, g}$ is surjective.
\end{definition}

\begin{definition}
  \label{def:H2}
  \lean{H2}
  \leanok
  \uses{def:dualnumber_C, lem:has_pullbacks}
  $(H_2)$ : if $Z = k$ and $Y = k[\varepsilon]$ then $p_{f, g}$ is bijective.
\end{definition}

\begin{definition}
  \label{def:H3}
  \lean{H3}
  \leanok
  \uses{lem:vsh2}
  $(H_3)$ : assuming $(H_2)$, $\dim_k (t_F) < \infty$.
\end{definition}

\begin{definition}
  \label{def:H4}
  \lean{H4}
  \leanok
  \uses{def:IsSmallExtension, lem:has_pullbacks}
  $(H_4)$ : if $f$ is a small extension then $p_{f, f}$ is bijective.
\end{definition}

These conditions are on the functor $F$, so for example "$F$ satisfies
$(H_1)$" means that for any morphisms $f$ and $g$, if $g$ is a
small extension, then $p_{f, g}$ is surjective.

\begin{lemma}
  \label{lem:vsh2}
  \leanok
  \lean{instPreservesLimitsOfShapeFGModuleCatResidueFieldDiscreteWalkingPairCompArtinResAlgTsze_functor}
  \uses{def:H2, lem:vstsze}
  The definition of $(H_3)$ makes sense because the condition $(H_2)$ implies
  that $F$ verifies the hypothesis of lemma~\ref{lem:vstsze}, so that
  $t_F$ admits a canonical $k$-vector space structure.
\end{lemma}

\begin{proof}
  \leanok
  \uses{lem:h2+}
\end{proof}

In fact, we have the
following result:

\begin{lemma}
  \label{lem:h2+}
  \lean{instPreservesLimitArtinResAlgDiscreteWalkingPairPairObjFGModuleCatResidueFieldTsze_functorOfH2}
  \leanok
  \uses{def:H2}
  If $F$ satisfies $(H_2)$, we have
  $F(A \times k[V]) \cong F(A) \times F(k[V])$
  for any $A \in \CL$ and any $V \in \Mod_k^{fg}$.
\end{lemma}

\begin{proof}
  \leanok
  By induction on the dimension of $V$.
\end{proof}

We also have the following generalization from $(H_1)$:

\begin{lemma}
  \label{lem:h1+}
  \lean{pullbackComparison_surj_of_surj}
  \leanok
  \uses{def:H1}
  If $F$ satisfies $(H_1)$, then $p_{f, g}$ is surjective
  for any surjection $g$.
\end{lemma}

\begin{proof}
  \leanok
  A surjection $g$ is either an isomorphism or it can be factored
  as a composition of small extensions.
  The result is clear for isomorphisms and we can show that
  if $g_1$ and $g_2$ are two suitable morphisms such that $p_{f, g_1}$
  and $p_{f, g_2}$ are surjective then $p_{f, g_2 \circ g_1}$
  is surjective as well.
\end{proof}

We will also need the following fact:

\begin{lemma}
  \label{lem:hRconds}
  \lean{instH1HR, instH2HR, instH3HR, instH4HR}
  \leanok
  \uses{def:hR, def:H1, def:H2, def:H3, def:H4, lem:H3hR}
  The functor $h_R$ satisfies conditions $(H_1)$, $(H_2)$,
  $(H_3)$ and $(H_4)$.
\end{lemma}

\begin{proof}
  \uses{lem:H3hR, lem:preservesPullbackshR}
  In fact in this case $(H_1)$, $(H_2)$ and $(H_4)$ hold trivially since
  $p_{f, g}$ is an isomorphism for any morphisms $f$
  and $g$, because $h_R$ is a hom-functor, and pullbacks in $\CL$ are
  pullbacks in $\Chat$.
  For $(H_3)$, it suffices to show the following linear isomorphism
  $t_R^* \cong \m_R / (\m_R ^ 2 + \m_\Lambda R)$.
\end{proof}

\chapter{Schlessinger's theorem}
\label{sec:thm}

Let $F$ be a functor from $\CL$ to $Sets$ such that $F(k)$ is a singleton.

\begin{theorem}
  \label{thm:Schlessinger1}
  \lean{Schlessinger1}
  \leanok
  \uses{def:H1, def:H2, def:H3, def:IsHull}
  If $F$ satisfies $(H_1)$, $(H_2)$ and $(H_3)$, then $F$ admits a hull.
\end{theorem}

\begin{theorem}
  \label{thm:Schlessinger2}
  \lean{Schlessinger2}
  \leanok
  \uses{def:H1, def:H2, def:H3, def:IsHull}
  If $F$ admits a hull, then $F$ satisfies $(H_1)$, $(H_2)$ and $(H_3)$.
\end{theorem}

\begin{theorem}
  \label{thm:Schlessinger3}
  \lean{Schlessinger3}
  \leanok
  \uses{def:H1, def:H2, def:H3, def:H4, def:Prorepresents}
  If $F$ satisfies $(H_1)$, $(H_2)$, $(H_3)$
  and $(H_4)$, then $F$ is pro-representable.
\end{theorem}

\begin{theorem}
  \label{thm:Schlessinger4}
  \lean{Schlessinger4}
  \leanok
  \uses{def:H1, def:H2, def:H3, def:H4, def:Prorepresents}
  If $F$ is pro-representable, then $F$ satisfies $(H_1)$, $(H_2)$, $(H_3)$
  and $(H_4)$.
\end{theorem}

\begin{proof}
  \leanok
  \uses{lem:hRconds, lem:vsnatlin}
  $h_R$ satisfies all four conditions as noted in \cref{lem:hRconds},
  so this is trivial
  (using \cref{lem:vsnatlin} for $(H_3)$).
\end{proof}

The main difficulty is the backward direction of $(1)$.

\chapter{Lemmas}
\label{sec:lemmas}

We now give some additional lemmas needed for the proof.

\section{Essential morphisms}

\begin{definition}
  \label{def:IsEssential}
  \lean{IsEssential}
  \leanok
  \uses{def:ArtinResAlg}
  A surjective morphism $p : B \to A$ in $\CL$ is \emph{essential}
  if for any morphism $q : C \to B$ such that $q \circ p$ is surjective,
  $q$ is surjective.
\end{definition}

\begin{lemma}
  \label{lem:essential_proartin}
  \lean{proartin_hom_surj}
  \leanok
  \uses{def:IsEssential, def:ProArtinResAlg}
  Let $p : B \to A$ be essential and $q : C \to B$ be a morphism in
  $\Chat$ such that $q \circ p$ is surjective. Then $q$ is surjective.
\end{lemma}

\begin{proof}
  \leanok
  \uses{lem:factorizes}
  This is trivial considering that $q$ factorizes through
  $q_n : C/\m_C^n \to B$ for some $n$ which is a morphism in $\CL$
  (\cref{lem:factorizes}).
\end{proof}

\begin{lemma}
  \label{lem:14ii}
  \lean{lem14ii}
  \leanok
  \uses{def:IsSmallExtension, def:IsEssential}
  If $p : B \to A$ is a small extension, then $p$ is not essential
  if and only if $p$ admits a section, that is there exists a morphism
  $s : A \to B$ such that
  $p \circ s = \mathbb{1}_A$.
\end{lemma}

\begin{proof}
  \leanok
  If $p$ has a section $s$, then $p$ is not essential since $s$
  is not surjective as we have
  $\length_\Lambda (B) = \length_\Lambda (A) + 1$.

  Now, suppose that $p$ is not essential, and let $q : C \to B$ be such that
  $q \circ p$ is surjective but $q$ is not. Let $C'$ be the
  range of $q$. We have $\length_\Lambda (C') < \length_\Lambda (B)$
  since $q$ is not surjective so
  $\length_\Lambda (C') \le \length_\Lambda (A)$,
  and $\length_\Lambda (C') \ge \length_\Lambda (A)$ because
  $p\restr{C'}$ is surjective (because $q \circ p$ is).
  Thus $p\restr{C'}$ is bijective
  which gives the section that we wanted.
\end{proof}

\section{Action of a small extension}

Let $p : B \to A$ be a small extension in $\CL$, and write $I = \ker p$.

\begin{lemma}
  \label{lem:isosmallext}
  \lean{rk215.iso}
  \leanok
  \uses{def:IsSmallExtension, def:tsze_C, lem:has_pullbacks}
  There is an isomorphism $B \times k[I] \cong B \times_A B$ whose
  first component is the identity.
\end{lemma}

\begin{proof}
  \leanok
  It is given by $(x, r + y) \mapsto (x, x + y)$.
  The inverse map is $(x, y) \mapsto (x, x_0 + y - x)$ where $x_0$ is
  the $k$ residue of $x$.
  It is easy to check that this defines an isomorphism.
\end{proof}

Using this isomorphism, we can get the following result
as in \cite[Remark 2.15]{Sch}:
\begin{lemma}
  \label{lem:action}
  \lean{rk215.action}
  \leanok
  \uses{lem:isosmallext, def:H2}
  Let $F : \CL \to Sets$ such that $F(k)$ is a singleton and $F$
  satisfies
  $(H_2)$. Using the isomorphism from lemma~\ref{lem:isosmallext}
  together with the comparison maps (\cref{eq:pfg}),
  we get a map
  $$F(B) \times F(k[I]) \to F(B) \times_{F(A)} F(B)$$
  This defines an action of $F(k[I])$ on $F(B)$, with orbits contained
  in the fibers of $p$.
\end{lemma}

\begin{proof}
  \leanok
  This is simply a formal check. In practice we only need
  the fact that $0$ acts trivially, so there is less to check.
\end{proof}

\begin{lemma}
  \label{lem:actionH1}
  \lean{rk215.action_surj}
  \leanok
  \uses{lem:action, def:H1}
  The condition $(H_1)$ implies that the map in \cref{lem:action}
  is surjective, i.e.
  the action is transitive on fibers of $p$.
\end{lemma}

\begin{proof}
  \leanok
\end{proof}

\begin{lemma}
  \label{lem:actionH4}
  \lean{rk215.action_bij}
  \leanok
  \uses{lem:action, def:H4}
  The condition $(H_4)$ implies that the map in \cref{lem:action}
  is bijective, i.e.
  the fibers of $p$ are principal homogeneous spaces under $F(k[I])$.
\end{lemma}

\begin{proof}
  \leanok
\end{proof}

\begin{lemma}
  \label{lem:actionnattrans}
  \lean{rk215.appcomm}
  \leanok
  \uses{lem:action}
  The action defined in \cref{lem:action} is functorial:
  if $G$ is another functor with the necessary conditions,
  and $\eta : F \to G$
  is a natural transformation, then the following diagram commutes:

\begin{center}
  % https://tikzcd.yichuanshen.de/#N4Igdg9gJgpgziAXAbVABwnAlgFyxMJZABgBoBGAXVJADcBDAGwFcYkQAxACgCEBKAAQAdIXgC28AdwDWyAJKU+IAL6l0mXPkIoATBWp0mrdt37DRWCXAD6wbgEE+yqbyWr12PASJliBhixsiCAA4q7m4pJhsgpuaiAYnlpEen40AcbBYWYikTbAYY7O2W4GMFAA5vBEoABmAE4QYkh6IDgQSADM7iANTUhkbR2I5D19zYiD7S1jjRPkNNOI3ZTKQA
  \begin{tikzcd}
    {G(B) \times G(k[I])} \arrow[rr]           &  & G(B) \times_{G(A)} G(B)           \\
    {F(B) \times F(k[I])} \arrow[rr] \arrow[u] &  & F(B) \times_{F(A)} F(B) \arrow[u]
  \end{tikzcd}
\end{center}
\end{lemma}

\begin{proof}
  \leanok
\end{proof}

\section{Topology}

We need the following lemma
\cite[\href{https://stacks.math.columbia.edu/tag/06SE}{Tag 06SE}]{stacks-project}:
\begin{lemma}
  \label{lem:topo}
  \lean{sametopo}
  \leanok
  \uses{def:complete}
  Let $R$ be an object of $\Chat$, and let $(J_n)_{n\in\N}$ be a decreasing
  sequence of ideals of $R$ such that $\m_R^n \le J_n$.
  Set $J = \bigcap_{n\in\N} J_n$.
  Then the sequence $(J_n/J)_{n\in\N}$ defines the $\m_{R/J}$-adic
  topology on $R/J$.
\end{lemma}

\begin{proof}
  \leanok
  We need to show that for any $n \in \N$, there exists $N\in\N$ such
  that $J_N \le J + \m_R^n$. For every $k \in \N$, $R / \m_R^k$
  is Artinian, so the decreasing sequence $(J_i + \m_R^k)_{i \in \N}$
  stabilizes. We write $N_k$ for some index such that $J_{N_k} + \m_R^k$
  is the minimum of the sequence.
  Note that we have
  $J_{N_k} + \m_R^k \le J_k + \m_R^k$ so since we also have
  $\m_R^k \le J_k$ we get that $J_{N_k} \le J_k$.
  We claim that $J_{N_n} \le J + \m_R^n$.
  Let $x \in J_{N_n}$. We define a Cauchy sequence $(x_i)$ in $R$ with
  $x_i \in J_{N_i}$ for $i \ge n$
  by induction, by first setting $x_i = x$ if $i \le n$.
  If $i > n$, then there exists
  $x_{i+1} \in J_{N_{i+1}}$ such that $x_i - x_{i+1} \in \m_R ^ i$:
  indeed
  if $N_{i+1} \le N_i$ we can take $x_{i+1}= x_i$, and otherwise we have
  $J_{N_{i+1}} + \m_R^i = J_{N_i} + \m_R^i$ (note that have $x_i \in J_{N_i}$).
  Since $R$ is complete, this sequence admits a limit $y$.
  We have $y \in J_{N_i} + \m_R^i$ for all $i$ so $y \in J$ and
  $y - x = y - x_n \in \m_R ^ n$, so $x \in J + \m_R^n$.
\end{proof}

In particular, this means that if $J \ne R$, we have
$\hat{F}(R/J) \cong \varprojlim F(R / J_n)$.

\chapter{Proof}
\label{sec:proof}

In this section we prove the main theorem (\cref{sec:thm}).

\section{Existence of a hull}
\label{mainproof}

The main proof is the following, showing the existence of a hull given
$(H_1)$, $(H_2)$ and $(H_3)$.
%TODO: split this
\begin{proof}
\proves{thm:Schlessinger1}
\leanok
\uses{lem:topo, lem:h1+, lem:smoothsmallext, lem:actionH1, lem:actionnattrans, lem:vsnatiso,
  lem:14ii, lem:essential_proartin, lem:powerseriesnoetherian, lem:mem_maximalideal_pow,
  lem:linearequiv_map, lem:quot_inf_equiv, lem:iscompletequot,
  lem:nattrans_of_comp_artinhom, lem:completion_map_quot_mk_eq_pi}

Suppose that $F$ satisfies $(H_1)$, $(H_2)$ and $(H_3)$.
Let $r = \dim_k(t_F)$, $S = \Lambda[[X_1, \ldots, X_r]]$ and $\n$
the maximal ideal of $S$.

We first construct $R$ as a projective limit
of quotients of $S$. More specifically we define a sequence
$(J_q, \xi_q)_{q\ge2}$ by induction where $(J_q)$ is a decreasing
sequence of ideals of $S$ such that $R_q = S / J_q \in \CL$ and
$\xi_q \in F(R_q)$ is such that $(\xi_q) \in \varprojlim F(R_q)$.

First we set $J_2 = \n^2 + \m_\Lambda S$ and $R_2 = S/J_2$.

There exists $\xi_2 \in R_2$ inducing an isomorphism
$\Hom (R_2, k[\varepsilon]) \cong t_F$.
To see this, first note that $R_2 \cong k[k^r] \cong k[\varepsilon]^r$.
By applying $(H_2)$, this yields $F(R_2) \cong t_F^r$,
so choosing a basis $(e_1, \ldots, e_r)$
of $t_F$ gives
an element $\xi_2 \in F(R_2)$. We then check that $u \mapsto F(u)(\xi_2)$
is an isomorphism. In fact, this map is linear and sends the morphism
$R_2 \to k[\varepsilon]$ corresponding to the $i$-th projection
$k^r \to k$
to $e_i$.

For the induction step, suppose that we have an ideal $J_q$ such
that $R_q = S/J_q \in \CL$ and $\xi_q \in F(R_q)$.
Consider the set $\mathscr{S}$ of ideals $J$ of $S$ such that
$\n J_q \le J \le J_q$ and $\xi_q$ lifts to $F(S/J)$.
%TODO: mention that S/J in \CL here
We will set $J_{q+1}$ to be the minimum of $\mathscr{S}$,
with $\xi_{q+1}$ some lift of $\xi_q$,
after proving that it exists.
First note that $\mathscr{S}$ is not empty as it contains $J_q$,
and since $S / \n J_q$ is Artinian, we know that
$\mathscr{S}$ admits a minimal element, so we only need
to show that $\mathscr{S}$ is stable by pairwise intersection
to get a minimum.
Let $J$ and $K$ be two elements of $\mathscr{S}$.
The elements of $\mathscr{S}$ are in correspondence with
$k$-vector subspaces of $J/\n J_q$,
so we can extend
$J$ into some ideal $J' \le J_q$ such that
$J' \cap K = J \cap K$ and $J' + K = J_q$, and we have
$J' \in \mathscr{S}$ because $J \le J'$.
We now use the isomorphism
$$S/J' \times_{S/J_q} S/K \cong S/(J \cap K)$$
together with \cref{lem:h1+} to conclude that $\xi_q$ lifts to
$S/(J \cap K)$, i.e. that $J \cap K \in \mathscr{S}$, which
concludes the induction step.


We set $J = \bigcap_{q \ge 2} J_q$ and $R = S / J$. By definition
of the $J_q$'s, we have $\n ^ q \le J_q$, so \cref{lem:topo} applies
and $\hat{F}(R) = \varprojlim F(S/J_q)$, so we can set
$\xi = \varprojlim \xi_q \in \hat{F}(R)$.

We have $\Hom (R, k[\varepsilon]) \cong \Hom (R_2, k[\varepsilon])$
because any map $S \to k[\varepsilon]$ factorizes through $S/J_2$,
and since $\xi$ projects to $\xi_2$ we have $t_R \cong t_F$
by definition of $\xi_2$.

It remains to show that $\nattrans(\xi) : h_R \to F$ is smooth.

Using \cref{lem:smoothsmallext}, let $p : A' \to A$ be a small extension,
$u : R \to A$, $\eta' \in F(A')$ and $\eta \in F(A)$ such that
$\nattrans(\xi) (u) = F(p)(\eta') = \eta$.
We need to find $u' : R \to A'$ such that $\nattrans(\xi) (u') = \eta'$
and $p \circ u' = u$.

First we reduce this to finding $u'' : R \to A'$ such that
$p \circ u'' = u$. Indeed, if we have such a $u''$, then
$\nattrans(\xi) (u'')$ and $\eta'$ both lie over $\eta$,
so by \cref{lem:actionH1}, there exists $\sigma \in F(k[I])$
sending $\nattrans(\xi) (u'')$ to $\eta'$,
where $I = \ker p$.
By using \cref{lem:actionnattrans} and the fact that
$h_R (k[I]) \cong F(k[I])$ (using \cref{lem:vsnatiso}),
we can also apply $\sigma$ to $u''$ giving us the $u' : R \to A'$
that we wanted.

Since $u$ factorizes through $R/\m_R^n \to A$ for some $n$,
and the topology is induced by the $J_q$'s (\cref{lem:topo}),
$u$ factorizes through $R_q \to A$ for some $q$.
It suffices to find $v$ completing the following diagram:
\begin{center}
\label{eq:proofdiagram}
% https://tikzcd.yichuanshen.de/#N4Igdg9gJgpgziAXAbVABwnAlgFyxMJZABgBpiBdUkANwEMAbAVxiRAGUQBfU9TXfIRQAmclVqMWbAEoB9AI4ACADrK8AW3iyAgou0BybrxAZseAkVEBGcfWatEIOfKN8zgomRvU7Ux3OB5AGorLm5xGCgAc3giUAAzACcIdSQyEBwIJCsfSQcQAHdXECSUtOpMpABmHgTk1MQqiqzEYVqS+uzmpFEJezY0RNkrYtKGpoyWnJAGLDB8qDo4AAtIkFz+xxpwriA
\begin{tikzcd}
  S \arrow[rr, "w"] \arrow[d]                 &  & R_q \times_A A' \arrow[d, "pr_1"] \\
  R_{q+1} \arrow[rr] \arrow[rru, "v", dashed] &  & R_q
\end{tikzcd}
\end{center}
where $w$ is defined using the universal property of
$S = \Lambda [[X_1, \ldots, X_r]]$ to make the diagram commute.
More specifically, to define $w$ we use that $pr_1$ is surjective so that
we can first define $w$ on the $X_i$'s, and since all maps are local the images
of the $X_i$'s are in the maximal ideal so we can extend the map to $S$.
In this context we can easily show that $pr_1$ is a small extension,
so we can apply \cref{lem:14ii}.
If $pr_1$ admits a section then we are done. Otherwise,
$pr_1$ is essential, and $pr_1 \circ w$ is the quotient map $S \to R_q$,
so $w$ is surjective by \cref{lem:essential_proartin}.
By $(H_1)$, $\xi_q$ lifts to $R_q \times_A A'$, so $\ker w \in \mathscr{S}$
and by minimality we have $J_{q+1} \le \ker w$. Thus $w$ factors through
$R_{q+1}$ giving the $v$ that we wanted.

\end{proof}

\section{Necessity of the conditions}
\label{necessityoftheconditions}

We now prove that the conditions $(H_1)$, $(H_2)$ and $(H_3)$
are necessary for the existence of a hull.

\begin{proof}
  \proves{thm:Schlessinger2}
  %TODO: split things
  \uses{lem:vsnatlin, lem:surj_of_smooth}
  Suppose that $F$ admits a hull $(R, \xi)$.

  First since $t_R \cong t_F$ (linearly by \cref{lem:vsnatlin}),
  and $h_R$ satisfies $(H_3)$, $F$ satisfies $(H_3)$.

  To check $(H_1)$, let $f : X \to Z$ and $g : Y \to Z$ be morphisms
  in $\CL$ with $g$ surjective, and $x \in F(X)$ and $y \in F(Y)$
  both lying over $z \in F(Z)$. First, since $h_R \to F$ is smooth,
  by \cref{lem:surj_of_smooth} there exists $u_x : R \to X$ such that
  $\nattrans(\xi) (u_x) = x$. Then by smoothness applied to $g$,
  there exists $u_y : R \to Y$ such that $\nattrans(\xi) (u_y) = y$
  and $g \circ u_y = f \circ u_x$.
  Then $\nattrans(\xi)(u_x \times u_y) \in F(X \times_Z Y)$
  projects to $x$ and $y$, so $p_{f, g}$ is surjective.

  To check $(H_2)$, we now only need to check injectivity, since
  we have already shown $(H_1)$.
  Let $A \in \CL$ and $\zeta_1$ and $\zeta_2$ in
  $F(A \times k[\varepsilon])$ having the same projections $a \in F(A)$
  and $e \in k[\varepsilon]$.
  We proceed similarly as in the proof of $(H_1)$.
  First there is $u' : R \to A$ such that $\nattrans(\xi) (u') = a$.
  Then by smoothness applied to the projection
  $pr_1 : A \times k[\varepsilon] \to A$, we get maps
  $u_i : R \to A \times k[\varepsilon]$ for $i = 1, 2$ such that
  $\nattrans(\xi) (u_i) = \zeta_i$ and $pr_1 \circ u_i = u'$.
  Since $\nattrans(\xi) (pr_2 \circ u_i) = e$ in both cases,
  and $\nattrans(\xi)_{k[\varepsilon]}$ is an isomorphism,
  we have $u_1 = u_2$ so $\zeta_1 = \zeta_2$.

\end{proof}

\section{Pro-representability}

We finally prove that the conditions $(H_1)$, $(H_2)$, $(H_3)$
and $(H_4)$ are sufficient for pro-representability.

\begin{proof}
  \proves{thm:Schlessinger3}
  \leanok
  \uses{thm:Schlessinger1, lem:surj_of_smooth, lem:smallext_obj_induction, lem:actionH4, lem:actionnattrans}
  Suppose that $F$ satisfies $(H_1)$, $(H_2)$, $(H_3)$ and $(H_4)$.
  By the first point of the theorem, we know that $F$ admits a hull
  $(R, \xi)$. Furthermore, since $\nattrans(\xi)$ is smooth, we
  already know
  that $\nattrans(\xi)_A$ is surjective for all $A \in \CL$
  (\cref{lem:surj_of_smooth}).
  We will show by induction on $\length_\Lambda (A)$ that it is
  injective, using \cref{lem:smallextexist}.
  If $A = k$, then this is true by definition of a hull.
  Otherwise, let $p : B \to A$ be a small extension, with $I = \ker p$,
  and suppose that $\nattrans(\xi)_A$ is an isomorphism.
  By applying \cref{lem:actionH4} together with \cref{lem:actionnattrans},
  it is easy check that $\nattrans(\xi)_B$ is injective. Namely we take
  $u_1$ and $u_2$ in $h_R(B)$ such that
  $\nattrans(\xi) (u_1) = \nattrans(\xi) (u_2)$, and we need to show that
  the action of $h_R(k[I])$ by $0$ sends $u_1$ to $u_2$,
  which is clear from the commutative diagram in \cref{lem:actionnattrans}.
\end{proof}

\chapter{Remaining \texttt{sorry}'s}
\label{sec:sorries}

Here is the list of all the remaining Lean statements using \texttt{sorry}
at the time of writing, with a mathematical description and some ideas
for the proofs.

\cref{lem:vsnatiso}

\begin{lemma}
  \label{lem:powerseriesnoetherian}
  \lean{instIsNoetherianRingPowerSeries_schlessinger}
  \leanok
  Let $R$ be a Noetherian ring. Then $R[[X]]$ is Noetherian.
\end{lemma}
This is necessary for \cref{mainproof} especially to show that some quotients
of $S = \Lambda[[X_1, \ldots, X_r]]$ are in $\CL$.
This is an important result that is currently missing from Mathlib. I
have heard that someone was working on this, so I shouldn't need to
prove this.


\begin{lemma}
  \label{lem:mem_maximalideal_pow}
  \lean{mem_maximalideal_pow}
  \leanok
  Let $R$ be a local ring, $\m$ its maximal ideal,
  $r, k \in \N$, $\n$ the maximal ideal of $R[[X_1, \cdots , X_r]]$
  and $x = (x_n)_{n \in \N ^ {r}} \in R[[X_1, \cdots , X_r]]$.
  Then $x \in \n ^ k \iff
  \forall n \in \N ^ {r}, x_n \in \m ^ {k - \deg (n)}$.
\end{lemma}
\begin{proof}
  The forward direction is easy.
  The converse should probably be shown by induction on $r$ (and $k$).
\end{proof}
This lemma is very important to deal with $S = \Lambda[[X_1, \ldots, X_r]]$,
for example to show that $S$ is complete when $\Lambda$ is
and to show that $R_2 \cong k[k^r]$ in \cref{mainproof}.



\begin{lemma}
  \label{lem:preservesPullbackshR}
  \lean{instPreservesLimitArtinResAlgWalkingCospanCospanHR}
  \leanok
  \uses{def:hR}
  The functor $h_R$ (see \cref{def:hR}) preserves pullbacks.
\end{lemma}
See \cref{lem:hRconds}.
\begin{proof}
  It suffices to show that the inclusion functor preserves pullbacks.
\end{proof}


\begin{lemma}
  \label{lem:H3hR}
  \lean{instH3HR}
  \leanok
  \uses{def:hR, def:H3}
  The functor $h_R$ satisfies $(H_3)$.
\end{lemma}
See \cref{lem:hRconds}.
\begin{proof}
  I have shown that $h_R(k[\varepsilon]) = Hom(R, k[\varepsilon])
  \cong Der_\Lambda (R, k[\varepsilon])
  \cong (\m_R / (\m_R ^ 2 + \m_\Lambda R))^*$.
  The main missing part is the linearity of the first bijection. The fact that
  $\m_R / (\m_R ^ 2 + \m_\Lambda R)$ is finite
  dimensional is also missing for now.
\end{proof}



\begin{lemma}
  \label{lem:linearequiv_map}
  \lean{proof1.linearequiv_map}
  \leanok
  \uses{def:tsze_functor, lem:vstsze}
  The map $Hom(k^r, k) \to Hom(k[k^r], k[\varepsilon])$ is $k$-linear ($r \in \N$).
  ($k$ here is the residue field of $\Lambda$)
\end{lemma}
The vector space structure on $Hom(k[k^r], k[\varepsilon])$ here is the one induced by
the coyoneda functor $Hom(k[k^r], -)$ (\cref{lem:vstsze}). This is somewhat similar
to the previous point, where the abstract vector space structure is the main issue.
This is used in \cref{mainproof} to show the existence of $\xi_2$.


\begin{lemma}
  \label{lem:quot_inf_equiv}
  \lean{proof1.quot_inf_equiv}
  \leanok
  Let $R$ be a $Λ$-algebra, and $J, K$ ideals of $R$. Then,
  $$R / J \times_{R/(J + K)} R / K \cong R / (J \cap K)$$
\end{lemma}
This is the easiest \texttt{sorry}.

\begin{lemma}
  \label{lem:iscompletequot}
  \lean{proof1.instIsCompleteQuotientMvPowerSeriesFinIdealIdeal}
  \leanok
  \uses{def:complete}
  In the proof of the theorem, we have $r \in \N$, $S = \Lambda[[X_1, \cdots, X_r]]$, and
  $J \neq S$ an ideal of $S$. $S / J$ is complete (as a local ring).
\end{lemma}
This is needed for \cref{mainproof}.
It is essentially a topological result that is in Mathlib
(as \lean{QuotientAddGroup.completeSpace'}),
however it is stated here in a way that doesn't mention topology,
and in fact the topology
has not been defined at all in Lean yet.

\cref{thm:Schlessinger2}
Only $(H_3)$ has been shown for now.
The proof is described in \cref{necessityoftheconditions}.

